\chapter{Reproducibility}
\label{app:repro}

This appendix records what would be needed to reproduce the results: the
configuration of every run family, and the provenance of every headline number.
The datasets themselves are described in Table~\ref{tab:datasets}, with counts
measured on the files rather than quoted.

\section{Hyperparameters}

% T12 -- Hyperparameters, per run family.
\begin{minipage}{\textwidth}
\centering\footnotesize
\captionof{table}{The training settings of each family of runs}
\label{tab:hyperparams}
\begin{tabular}{@{}l@{\hspace{6pt}}>{\raggedright\arraybackslash}p{2.6cm}@{\hspace{6pt}}>{\raggedright\arraybackslash}p{2.9cm}@{\hspace{6pt}}>{\raggedright\arraybackslash}p{2.9cm}@{\hspace{6pt}}>{\raggedright\arraybackslash}p{3.1cm}@{}}
\toprule
& \textbf{Iter.\ 1} & \textbf{Iter.\ 2--6} & \textbf{Iter.\ 7--9} & \textbf{Iter.\ 10--11} \\
\midrule
Base model        & Mistral-7B-v0.1 & Qwen2.5-0.5B-Instruct & Qwen2.5-0.5B-Instruct & Qwen2.5-0.5B-Instruct \\
Loading           & 4-bit (QLoRA)   & bf16 / fp32           & bf16 / fp32           & bf16 / fp32 \\
LoRA rank $r$     & 16              & 16                    & 16                    & 16 \\
LoRA $\alpha$     & 32              & 32                    & 32                    & 32 \\
LoRA dropout      &                 & 0.1                   & 0.1                   & 0.1 \\
Learning rate     & default         & $1.5\times10^{-4}$    & $1.5\times10^{-4}$    & $1.5\times10^{-4}$ \\
Weight decay      &                 & 0.01                  & 0.01                  & 0.01 \\
Max passes        & 100 steps       & 20                    & 20                    & 20 (cosine horizon) \\
Early stopping    &                 & patience 8, then 4    & patience 4            & 4 evaluations \\
Eval.\ sequence length &            & 512 tokens            & 512 tokens            & 512 tokens \\
Training window   &                 & 128 events, stride 64 & 128 events, stride 64 & 128 events, stride 64 \\
Context fraction  &                 & 0.8                   & 0.8                   & 0.8 \\
Added symbols     & antenna names   & 369 antennas          & 369 antennas + 24 hours & 369 antennas + 24 hours \\
Transition sampler &                &  $\times3$ from iter.\ 6 & $\times3$             & $\times3$ \\
Emphasised hours  &                 & global (4--6, 18--20) & per kind of user      & per kind of user \\
Kinds of user     &                 &                       & 4, read from the past & 4, refitted at each size \\
Stopping protocol & fixed steps     & passes                & passes                & 1{,}000 optimiser steps \\
Seed              &                 & 42                    & 42                    & 42 \\
Hardware          & Colab T4        & Colab T4              & Colab T4              & Colab T4, 42 min/size \\
\bottomrule
\end{tabular}

\end{minipage}


\section{Provenance of the headline numbers}

% T13 -- Which artefact produces which headline number.
\begin{minipage}{\textwidth}
\centering\footnotesize
\captionof{table}{The origin of each number quoted in the report}
\label{tab:provenance}
\begin{tabular}{@{}>{\raggedright\arraybackslash}p{4.5cm}ll>{\raggedright\arraybackslash}p{6.0cm}@{}}
\toprule
\textbf{Number quoted in the report} & \textbf{Data} & \textbf{Iter.} & \textbf{Which run produced it} \\
\midrule
16.87\,\% first signal              & \sraw & 1  & single-notebook Mistral-7B QLoRA run, 100 steps \\
48.7\,\% $\rightarrow$ 64.4\,\%     & \sraw & 2  & same notebook, keeping the best saved copy \\
$r=-0.971$ on the switch rate       & \textsc{raw day}, 500 people & 2 & per-person analysis of the 1{,}173-person grading \\
$+1.1$ at the 80/20 cut             & \sfilt & 4  & training run on the filtered day \\
No cut beats the rule               & \smob & 5  & one trained model, seven grading passes \\
$-3.0$ from shuffling the hours     & \smob & 6  & re-grading the same trained copy \\
$+2.6$ at 18--20h                   & \smob & 6  & the targeted-sampling run, whole day graded \\
$27.4$ separation of the derived hours & \smob & 7 & analysis of the kinds of user on the training file \\
$+3.7$, first run above the rule    & \smob & 8  & routing each kind of user to its own numbers \\
$+4.2$ at 1{,}900 training people   & \stwok & 9  & scaling run on the 2{,}000-person day \\
$-1.6$ on stretched days            & \slev & 9  & the 200-record levelling run \\
$+4.3$, flat from 5{,}000 to 21{,}000 & \seleven & 10 & four-size scaling notebook, constant training time \\
The resampling intervals            & \seleven & 10 & 2{,}000 paired resamples over the test people \\
0 of 7 weekday trainings finished   & \seleven & 11 & seven-day notebook; the session crashed after 5\,h \\
\bottomrule
\end{tabular}

\figtext{%
\figpara{Table~\ref{tab:provenance}: data provenance mapping for experimental results.}{Every number quoted anywhere in the body can be
traced back to one run through Table~\ref{tab:provenance}. The \textbf{Data} column names the set
of people it was measured on, described in \textbf{Table~\ref{tab:datasets}}, and
\textbf{Iter.} is the numbered iteration of \textbf{Appendix~\ref{app:timeline}}.}

\figpara{Table~\ref{tab:provenance}: reproducibility tracking across notebook runs.}{Two rows in Table~\ref{tab:provenance} measured on
different sets of people are not comparable on their absolute scores, only on
their distances to the rule of that row's own data. Reading Table~\ref{tab:provenance} alongside
\textbf{Table~\ref{tab:datasets}} is what makes that check possible.}}

\end{minipage}


\section{Artefacts}

Three notebooks and a small data-preparation pipeline produce everything above.
Table~\ref{tab:artefacts} lists them factually; none of their internals is
described anywhere in this report, by design.

\begin{table}[t]
\centering\footnotesize
\caption[Artefacts]{The artefacts, and what each one produces.}
\label{tab:artefacts}
\begin{tabular}{@{}ll>{\raggedright\arraybackslash}p{6.2cm}@{}}
\toprule
\textbf{Artefact} & \textbf{Scope} & \textbf{Produces} \\
\midrule
\texttt{train\_cellid\_llm.ipynb}   & one training run   & iterations 1--8: all results up to $+3.7\pt$ \\
\texttt{train\_cellid\_final.ipynb} & four nested tiers  & iterations 10--11: the scaling ladder and its bootstrap \\
\texttt{train\_cellid\_week.ipynb}  & seven weekdays     & the day-of-week study, which did not complete \\
\texttt{make data-train}            & data preparation   & the train/test JSONL files, timestamps kept as hours \\
\texttt{make data-merge}            & data preparation   & \seleven: one merged training file, nested test sets \\
\texttt{make groups}                & analysis           & the behavioural-group report of Chapter~\ref{ch:groups} \\
\bottomrule
\end{tabular}
\end{table}

\section{Configuration keys that carry a methodological decision}

A handful of configuration values are not tuning knobs but encode a decision
argued in the body. They are listed here so that a reader who changes one knows
which claim they are changing.

\begin{table}[t]
\centering\footnotesize
\caption[Configuration keys carrying a decision]{Configuration values that
encode a methodological choice rather than a tuning preference.}
\label{tab:config-keys}
\begin{tabular}{@{}ll>{\raggedright\arraybackslash}p{6.7cm}@{}}
\toprule
\textbf{Key} & \textbf{Value} & \textbf{Decision it encodes} \\
\midrule
\texttt{context\_fraction}              & 0.8            & the reference cut; part of the ruler (Ch.~\ref{ch:first}) \\
\texttt{max\_epochs}                    & 20             & also the cosine horizon: lowering it changes the regime \\
\texttt{early\_stop\_patience}           & 4              & the validation peak is early on this task \\
\texttt{transition\_windows}             & (4,6,$\times2$), (18,20,$\times4$) & the \emph{guessed} global window; fallback only \\
\texttt{group\_windows}                  & \texttt{None}  & derive the windows from the data, per group (Ch.~\ref{ch:groups}) \\
\texttt{transition\_chunk\_oversample}   & 3.0            & oversample windows containing a transition \\
\texttt{group\_bands}                    & 4 bands        & coarse enough for each band to be estimable \\
\texttt{group\_shrinkage}                & 8.0            & a sparse band must not dominate a profile \\
\texttt{group\_detect\_fraction}         & prefix only    & no leakage: the group is read from the past \\
\texttt{max\_steps} (scaling)            & 1{,}000        & constant compute across tiers (Ch.~\ref{ch:scale}) \\
\texttt{val\_subsample} (scaling)        & 250 fixed users& the same stopping rule at every tier \\
\bottomrule
\end{tabular}
\end{table}
